\chapter*{Code of conduct and transparency statement on the use of GenAI for KU Leuven-students (academic year 2025--2026)}
\addcontentsline{toc}{chapter}{Code of Conduct}
\markboth{Code of Conduct}{Code of Conduct}
Generative AI (GenAI) assistance tools can be used to generate various types of content, including text, images, code, video, music, or combinations thereof. Common examples of such tools include ChatGPT, Google Gemini, Microsoft Copilot, Midjourney, Claude.ai, Perplexity.ai, and DALL·E, among others.
\newline
\newline
This code of conduct is a tool that helps students to be transparent about the use of GenAI and fits within the university’s principles on academic integrity.

\subsection*{Important guidelines and remarks}

\begin{itemize}
    \item \textbf{Sensitive or personal data:} Some GenAI tools protect your input and sensitive or personal data better than others. There is often no transparency on what the owners of the AI applications do with the data entered. Therefore, \textbf{do not enter sensitive or personal data in free GenAI tools.} More info about what sensitive and personal data are, is described in the \href{https://www.kuleuven.be/english/education/leuvenlearninglab/support/toolguide/guidelines-for-safe-use-of-genai-tools}{three confidentiality levels of the KU Leuven data classification model} (non-confidential, confidential, strictly confidential). For confidential data you should preferably use \textbf{M365 Copilot Chat with your KU Leuven account}. If you use another GenAI tool, you must be \textbf{absolutely certain} it does \textbf{not store or reuse the data you enter. Try to avoid entering these data. Do so only if strictly necessary, and only to the extent required. For personal data, work with anonymous or pseudonymized data.} Additionally, in case of strictly confidential data this should first be discussed with the teaching staff of the course or your thesis supervisor.
    
    \item \textbf{Copyrighted materials:} For \textbf{lawfully obtained copyrighted material} you should preferably use M\textbf{365 Copilot Chat with your KU Leuven account.} If you use another GenAI tool, you must be \textbf{absolutely certain} it does n\textbf{ot store or reuse the data you enter.}
    
    \item GenAI assistance may not be used for data or topics covered by a \href{https://www.kuleuven.be/english/stuvo/student-legal-services/research-with-interaction-with-a-company-or-organisation}{Non-Disclosure Agreement (NDA)}. Even in the absence of an NDA, certain information may still need to be treated as confidential—for example, due to regulatory requirements or the risk of significant harm to the university if disclosed. In case of doubt, check with your teaching staff or supervisor.
    
    \item If your master’s thesis is under \href{https://iiw.kuleuven.be/english/students/master-thesis/embargo}{embargo}, you should first discuss with your supervisor whether the use of GenAI is permitted.
    
    \item Before using a GenAI tool, always consider whether its use is responsible, including from a \href{https://www.kuleuven.be/english/genai/sustainable-and-responsible-use-of-genai}{sustainability perspective} (e.g. when using GenAI as a search engine, language assistant,…).
    
    \item \textbf{Take a scientific and critical attitude} when interacting with GenAI assistance and interpreting its output, that may not always be correct. 
    
    \item As a student you are responsible for complying with Article 84 of the Regulations on Education and Examinations: your report or thesis should reflect your own knowledge, understanding and skills. Be aware that plagiarism rules also apply to (work that is the result of) the use of GenAI assistance tools. 
    \newline
    \textbf{Exam Regulations Article 84:} \textit{“Every conduct individual students display with which they (partially) inhibit or attempt to inhibit a correct judgement of their own knowledge, understanding and/or skills or those of other students, is considered an irregularity which may result in a suitable penalty. A special type of irregularity is plagiarism, i.e. copying the work (ideas, texts, structures, designs, images, plans, codes , ...) of others or prior personal work in an exact or slightly modified way without adequately acknowledging the sources. Every possession of prohibited resources during an examination (see article 65) is considered an irregularity.”}
    
    \item In order to maintain academic integrity and avoid plagiarism, \textbf{more information about being transparent on the use of GenAI assistance and about correctly citing and referencing GenAI} can be found on this website for students (\href{https://www.kuleuven.be/onderwijs/student/onderwijstools/artificiele-intelligentie#Transparantie}{Dutch}/\href{https://www.kuleuven.be/english/education/student/educational-tools/generative-artificial-intelligence#Transparency}{English}). 
    
    \item \textbf{Additional reading: KU Leuven guidelines on responsible use of Generative AI tools, and other information }(\href{https://www.kuleuven.be/genai}{Dutch}/\href{https://www.kuleuven.be/english/genai}{English})
\end{itemize}

\subsection*{A few final words}

\textbf{If you are uncertain whether or not you should declare your use of GenAI tools, we suggest that you discuss this with your instructor or supervisor. It is always safer to declare GenAI use, even when it is not strictly required. However, declaring GenAI use does not entail that its use is allowed; the right column in the table below provides more detailed instructions in this regard (code of conduct). }
\newline
\newline
Moreover, advanced AI tools are evolving rapidly, and their capabilities have expanded significantly in a short period of time. As a result, we do not yet have all the answers about their responsible use. Finally, it is important to follow-up on the most recent evolutions in AI technologies, to have an open mind but also to be a bit cautious, to communicate with instructors, teaching assistants, supervisors and peers, to be as transparent as we can, and to learn together as we move along.
\newpage
% --- Student info row ---
% --TO BE FILLED IN--
\noindent
\textbf{Student name:} \dotfill \hfill \textbf{Student number:} \dotfill

\vspace{1em}

% --- Grey box section title ---
\noindent
\colorbox{lightgray}{\parbox{\dimexpr\linewidth-2\fboxsep\relax}{\textbf{Please indicate with "X" whether it relates to a course assignment or to the Bachelor’s or Master’s thesis:}}}

\vspace{1em}

O This form is related to a \textbf{course assignment.}

\begin{tabularx}{\linewidth}{@{}lXlX@{}}
\textbf{Course name:} & \dotfill & \textbf{Course number:} & \dotfill \\
\end{tabularx}

\vspace{1em}

This form is related to my O \textbf{Bachelor’s} or O \textbf{Master’s thesis.}

\begin{tabularx}{\linewidth}{@{}lXlX@{}}
\textbf{Title Bachelor’s or Master’s thesis:} & \dotfill & \textbf{Supervisor:} & \dotfill \\
\end{tabularx}

\vspace{1em}

\textbf{Daily supervisor:} \dotfill

\vspace{2em}

% --- Grey box section title ---
\noindent
\colorbox{lightgray}{\parbox{\dimexpr\linewidth-2\fboxsep\relax}{\textbf{Please indicate with "X":}}}

\vspace{1em}

O I \textbf{did not use} any GenAI assistance tool. \\[1em]

O I \textbf{did use} GenAI Assistance. In this case \textbf{specify which ones} (e.g.\ ChatGPT, M365 Copilot, …): \dotfill

\newpage
% ====== PREAMBLE AANVULLINGEN ======

\definecolor{lightgray}{gray}{0.9}

% Kolomdefinitie: mooie linkslijnende X-kolom
\newcolumntype{Y}{>{\raggedright\arraybackslash}X}

% Veilige cel-wrapper (geen \setlist hier!)
\newcommand{\Cell}[1]{\begin{minipage}[t]{\linewidth}\raggedright\small #1\end{minipage}}

% (optioneel) iets compactere tabel
\setlength{\tabcolsep}{6pt}
\renewcommand{\arraystretch}{1.2}

\begin{longtable}{|p{0.25\textwidth}|p{0.2\textwidth}|p{0.55\textwidth}|}
\hline 
\textbf{GenAI assistance used as/for:} & \textbf{Name of the GenAI tool(s) used} If helpful, also describe in which way you were using GenAI related to what is specified as code of conduct. & \textbf{Code of conduct: \newline
For each of the categories below, always take into account the important guidelines and remarks mentioned above (e.g. copyrighted data, sensitive or personal data,…).} \\ \hline
\endfirsthead

%\hline
%\textbf{GenAI assistance used as/for:} & \textbf{Name of the GenAI tool(s) used} If helpful, also describe in which way you were using GenAI related to what is specified as code of conduct. & \textbf{Code of conduct:} For each of the categories below, always take into account the important guidelines and remarks mentioned above (e.g. copyrighted data, sensitive or personal data,…). \\  \\ \hline
%\endhead

\hline
\endfoot

\hline
\endlastfoot

\textbf{As a language assistant for reviewing or improving texts I wrote myself} & \textit{add text here} & This use is similar to using spelling and grammar check tools. In general, you do not have to refer to such kind of GenAI use in the text. 

\textbf{However, be careful: }
\begin{itemize}
    \item When using GenAI tools on texts you did not write yourself to improve the text, \textbf{you have to \href{https://bib.kuleuven.be/english/training-and-tutorials/citation}{refer} to the original source or author}, otherwise you are committing \href{https://www.kuleuven.be/english/education/plagiarism/index}{plagiarism} and thus an irregularity.
\end{itemize} \\ \hline
\textbf{As a paraphrasing tool} & \textit{add text here} & This use is allowed except when it is prohibited by the teacher or the program of study. You may paraphrase your own text or texts by an author other than yourself and take inspiration from what a GenAI tool or another tool suggests (unless it is not allowed). In general, you do not have to refer to such kind of GenAI use or other paraphrasing tools in the text. 

\textbf{However, be careful}: 
\begin{itemize}
    \item If it entails text by an author other than yourself, \textbf{you are not allowed to include that paraphrased text without \href{https://bib.kuleuven.be/english/training-and-tutorials/citation}{reference} to the original source or author}. Without such reference, you would be committing \href{https://www.kuleuven.be/english/education/plagiarism/index}{plagiarism} and thus an irregularity.
\end{itemize} \\ \hline
\textbf{For translation aid to improve texts I wrote myself or to better understand text from others} & \textit{add text here} & This use is allowed except when it is prohibited by the teacher or the program of study. It is similar to using translation tools (Google translate, DeepL, …). In general, you do not have to refer to such kind of GenAI use in the text. 

\textbf{However, be careful:} 
\begin{itemize}
    \item \textbf{You are not allowed to include that translated text without \href{https://bib.kuleuven.be/english/training-and-tutorials/citation}{reference} to the original source or author.} Without such reference, you would be committing \href{https://www.kuleuven.be/english/education/plagiarism/index}{plagiarism} and thus an irregularity.
    \item Always check the translated text for correctness and meaning.
\end{itemize} \\ \hline
\textbf{As a search engine to get information on a topic or to search for existing research on the topic} & \textit{add text here} & This use is similar to e.g. a Google search or checking Wikipedia. If you write your own text based on this information, you do not have to refer to the use of GenAI in the text. \textbf{You only have to \href{https://bib.kuleuven.be/english/training-and-tutorials/citation}{refer} to the existing research and references you have checked and used }(without such references you would be committing \href{https://www.kuleuven.be/english/education/plagiarism/index}{plagiarism} and thus committing an irregularity).

\textbf{However, be careful:}
\begin{itemize}
    \item Be aware that the output of the GenAI tool cannot be guaranteed as a 100\% reliable source of information. The output may not be entirely correct and/or be limited due to the databases it uses. Moreover, knowledge evolves and may change over time; therefore, the database of the GenAI tool may not be up to date. Therefore, \textbf{verify the information and do not just copy-paste it} as you should understand and critically process everything you are writing.
\end{itemize} \\ \hline
\textbf{For literature search} & \textit{add text here} & This use is comparable to e.g.\ a Google Scholar search. You do not have to refer to such kind of GenAI use; \textbf{you only have to \href{https://bib.kuleuven.be/english/training-and-tutorials/citation}{refer} to the literature references you have checked and used} (without such references you would be committing \href{https://www.kuleuven.be/english/education/plagiarism/index}{plagiarism} and thus committing an irregularity).

\medskip
\textbf{However, be careful:}
\begin{itemize}
    \item Be aware that the search output is restricted to the database the GenAI tool is built on. After this initial search, look for scientific sources and conduct your own analysis of the source documents. \textbf{Interpret, analyse and process the information you obtained; verify it and do not just copy-paste it} as you should understand and critically process everything you are writing.
    \item Be aware that some GenAI tools may \textbf{output \href{https://www.kuleuven.be/english/education/leuvenlearninglab/support/toolguide/guidelines-for-safe-use-of-genai-tools}{no or wrong references.}} As a student you are responsible for further checking and verifying the absence or correctness of references; do not just copy-paste it.
\end{itemize} \\ \hline
\textbf{For generating programming code } & \textit{add text here} & Use of GenAI for coding is allowed except when it is prohibited by the teacher or the program of study. If used for coding, correctly mention the use of GenAI assistance in accordance with the instructions on the \href{https://www.kuleuven.be/english/education/student/educational-tools/generative-artificial-intelligence#Transparency}{page on being transparant about the use of GenAI.} \\ \hline
\textbf{For generating new (research) ideas} & \textit{add text here} & This use of GenAI is allowed except when it is prohibited by the teacher or the program of study. Correctly mention the use of GenAI assistance in accordance with the instructions on the \href{https://www.kuleuven.be/english/education/student/educational-tools/generative-artificial-intelligence#Transparency}{page on being transparent about the use of GenAI.}

\medskip
\textbf{Be careful:}
\begin{itemize}
    \item Further verify in this case whether the idea is novel or not. It is likely that it is related to existing work. \textbf{If so, that existing work should be correctly \href{https://bib.kuleuven.be/english/training-and-tutorials/citation}{referenced} in the text} (without such reference you would be committing \href{https://www.kuleuven.be/english/education/plagiarism/index}{plagiarism} and thus committing an irregularity).
\end{itemize}\\ \hline
\textbf{For generating synthetic data} & \textit{add text here} & Use of GenAI for generating synthetic data is allowed, \textbf{provided that it is methodologically and ethically justifiable}, except when it is prohibited by the teacher or the program of study. Always correctly mention the use of GenAI assistance in accordance with the instructions on the \href{https://www.kuleuven.be/english/education/student/educational-tools/generative-artificial-intelligence#Transparency}{page on being transparent about the use of GenAI.}

\medskip
\textbf{Be careful:}
\begin{itemize}
    \item Always \textbf{carefully evaluate} the generated synthetic data for quality and possible bias since the output is highly dependent on the quality of the data on which the models are trained.
\end{itemize} \\ \hline
\textbf{For generating blocks of text
(other than the allowed use without referencing mentioned above)
} &  & According to Article 84 of the Regulations on Education and Examinations your text should allow to correctly and properly assess your own knowledge, understanding and skills. \textbf{Therefore, inserting blocks of text without quotes and a reference to GenAI assistance in your work is not allowed. }

\medskip
\textbf{Be careful:}
\begin{itemize}
    \item If it is really needed to insert a block of text from a GenAI tool, for instance because of the nature of your assignment, mention it as a citation by using quotes and correctly mention the use of GenAI assistance in accordance with the instructions \href{https://www.kuleuven.be/english/education/student/educational-tools/generative-artificial-intelligence#Transparency}{page on being transparent about the use of GenAI.}
    \item However, in general, such GenAI use should be kept to an \textbf{absolute minimum}; you should always check the original sources.
\end{itemize} \\ \hline
\textbf{For generating visuals, video or audio} & \textit{add text here} & Use of GenAI for generating visuals, audio or video is allowed except when it is prohibited by the teacher or the program of study. 

\medskip
\textbf{Be careful:}
\begin{itemize}
    \item If used, refer to GenAI for visuals according to the usual referencing style, following the instructions on the \href{https://bib.kuleuven.be/english/training-and-tutorials/citation/refering-to-genai}{page on referencing GenAI.}
    \item If you work with existing visuals, audio or video, \textbf{that existing work should be correctly \href{https://bib.kuleuven.be/english/training-and-tutorials/citation}{referenced} in the text} (without such reference you would be committing \href{https://www.kuleuven.be/english/education/plagiarism/index}{plagiarism} and thus an irregularity).
    \item Explain the usage in the methods section (if there is one) and optionally attach (or link to) the prompts with the full output (or history).
\end{itemize} \\ \hline
\textbf{Other use (specify here; this may also include a combination of types of use mentioned above):} & \textit{add text here} & To make sure other use of GenAI is allowed within the course or thesis, and if so, the conditions that may apply, contact the teaching staff of the course or the supervisor of the thesis \textbf{beforehand} and explain the intended GenAI use. Also inform the programme director. 

\medskip
Motivate how you would comply with Article 84 of the exam regulations. Explain the use and the added value of the AI tool you consider to use and how it is in accordance with the assignment or thesis and the \href{https://www.kuleuven.be/english/genai}{KU Leuven guidelines on responsible use of Generative AI tools.} 

\medskip
Depending on the kind of GenAI use, it may be needed to properly reference it in the text, in accordance with the instructions \href{https://www.kuleuven.be/english/education/student/educational-tools/generative-artificial-intelligence#Transparency}{page on being transparent about the use of GenAI.}\\ \hline
\end{longtable}